\chapter{Poder estatístico}
\label{sec:poder}

O cálculo de poder responde a uma pergunta de planejamento: que efeito o desenho
é capaz de detectar? Compara-se o efeito mínimo detectável com o efeito que a
literatura autoriza esperar. O que segue é um cálculo com parâmetros supostos e
contagens públicas de nascimentos do pré-período; será refeito com os
microdados antes de qualquer estimação.

\section{O efeito que se espera}

O efeito de referência é o de \citeonline{calzada2023bananas}: redução de 80 gramas
no peso ao nascer, com erro-padrão de 31, para filhos de mães que moram a até
150 metros de bananal fumigado por avião e cuja gestação atravessa a estação de
fumigação intensa. Além de 150 metros os autores não encontram efeito.

O que a proibição produz na média de um município é esse efeito multiplicado
pela fração $f$ de nascimentos expostos. No Equador, 1,8\% dos nascimentos
geocodificados do país estão a até 150 metros de um bananal, em mais de 186 mil
hectares de plantações. Transposta essa densidade aos mil hectares da Chapada,
chega-se a algo como trinta nascimentos expostos por ano. Limoeiro do Norte e
Quixeré registraram em média 1.068 nascimentos por ano entre 2010 e 2018; com
Russas, 2.078.\footnote{IBGE, Estatísticas do Registro Civil, tabela SIDRA
2609: nascidos vivos registrados no ano, por lugar de residência da mãe, média
de 2010 a 2018, extraída em 16/11/2026. A mesma tabela dá 8.874 nascimentos
anuais para os vinte e um municípios de controle. O denominador definitivo será
o do SINASC.} A fração exposta fica entre 1,4\% e 2,8\%, e possivelmente abaixo
disso, porque a costa equatoriana é mais povoada em torno das plantações.

\begin{table}[ht]
\centering\small
\caption{Efeito esperado sobre o peso médio municipal ao nascer (gramas), no
cenário em que havia aplicação aérea até a lei}
\label{tab:esperado}
\begin{tabular}{@{}lrrrr@{}}
\toprule
 & \multicolumn{4}{c}{Fração de nascimentos expostos} \\
\cmidrule(l){2-5}
Efeito por exposta & 1,4\% & 2,8\% & 5\% & 10\% \\
\midrule
$-80$ g (estimativa pontual) & $-1{,}1$ & $-2{,}2$ & $-4{,}0$ & $-8{,}0$ \\
$-140$ g (limite do intervalo de 95\%) & $-2{,}0$ & $-3{,}9$ & $-7{,}0$ & $-14{,}0$ \\
\bottomrule
\end{tabular}
\fonte{Elaboração própria.}
\end{table}

O caso central é de um a dois gramas. Três ressalvas. A tabela vale para o
cenário em que o avião voava até a lei; se já não voava, o efeito esperado do
contraste de 2019 é zero. O efeito de referência é o de gestações que
atravessam a estação intensa, e a tabela o aplica a todos os nascimentos
próximos, o que a torna um teto. No sentido oposto, ele mede a intensificação
sazonal entre quem já mora ao lado do bananal, e uma proibição remove toda a
fumigação. A fração exposta é uma transposição, não uma medida; pode ser
substituída pelo cruzamento dos setores censitários com a área de banana.

\section{O efeito mínimo detectável}

Considere-se o contraste mais simples: a diferença, entre tratados e controles,
da variação do peso médio entre todo o pré-período e todo o pós-período. Para
nível de significância $\alpha$ e poder $1-\beta$,
\[
EMD = (z_{1-\alpha/2} + z_{1-\beta}) \times
\sqrt{\sigma^2\Big(\tfrac{1}{N_{T}^{\text{pré}}} + \tfrac{1}{N_{T}^{\text{pós}}} +
\tfrac{1}{N_{C}^{\text{pré}}} + \tfrac{1}{N_{C}^{\text{pós}}}\Big) + 2\,\tau^2\Big(\tfrac{1}{J_T^{*}}+\tfrac{1}{J_C^{*}}\Big)},
\]
em que $\sigma$ é o desvio-padrão individual do peso ao nascer; os $N$ são os
números de nascimentos; $\tau$ é o desvio-padrão do choque comum aos nascimentos
de um mesmo município em cada um dos dois blocos de tempo, suposto independente
entre blocos; e $J^{*}$ é o número efetivo de municípios de cada grupo, igual
ao inverso da soma dos quadrados dos pesos. O segundo termo é o que a
correlação dentro do município acrescenta. Com poucos municípios tratados, ele
não diminui com o número de nascimentos, e a ponderação por nascimentos reduz
$J_C^{*}$ abaixo de vinte e um, porque um dos controles é uma cidade grande.

A Tabela~\ref{tab:emd} dá o valor com $\tau = 0$, para $\alpha = 5\%$, poder de
80\% e $\sigma = 500$ gramas, com o pré-período de 2010 a 2018 e seis anos de
pós. É um limite inferior.

\begin{table}[ht]
\centering\small
\caption{Limite inferior do efeito mínimo detectável sobre o peso ao nascer
($\tau = 0$)}
\label{tab:emd}
\begin{tabular}{@{}lrrrr@{}}
\toprule
 & & & \multicolumn{2}{c}{Limite inferior do EMD (g)} \\
\cmidrule(l){4-5}
Grupo tratado & Municípios & Nascimentos/ano & Pré agrupado & Base de um ano \\
\midrule
Limoeiro do Norte e Quixeré & 2 & 1.068 & 23,9 & 49,0 \\
Com Russas                   & 3 & 2.078 & 18,0 & 36,9 \\
\bottomrule
\end{tabular}
\fonte{Elaboração própria. Considera só o erro amostral. Os valores de $\sigma$ e $\tau$ são hipóteses de planejamento. A última coluna compara os seis anos de pós com um único ano de base, que é como o estimador proposto constrói cada efeito; por período isolado, o valor é maior.}
\end{table}

Três observações sobre a tabela. O valor de 18 a 24 gramas vale para o contraste
agrupado; os efeitos por período do estudo de evento, comparados a uma base
curta, têm efeito mínimo detectável pelo menos duas vezes maior. O valor
verdadeiro é ainda maior, porque $\tau$ não é zero e há dois ou três municípios
tratados, muito abaixo do número de conglomerados que a prática de avaliação
recomenda por grupo. E incluir Russas reduz o efeito mínimo detectável ao custo
de dobrar os nascimentos sem acrescentar exposição medida: o efeito esperado
cai pela metade.

Ficam para a fase de pesquisa, com os microdados e só com o pré-período: o
desvio-padrão observado do peso ao nascer, a estimativa de $\tau$, o efeito
mínimo detectável do desenho de calendário e o dos desfechos binários e raros
(baixo peso, óbito fetal, razão de sexos, intoxicação).

\section{O que a comparação diz}

O efeito esperado sobre a média municipal é da ordem de dez vezes menor que o
limite inferior do efeito mínimo detectável. Para que o contraste principal
detectasse o efeito de referência, cerca de 30\% dos nascimentos teriam de ser
de mães expostas, o que não é plausível para mil hectares. Nenhuma das
especificações propostas tem, com as hipóteses atuais, poder para detectar o
efeito de referência na média municipal.

O projeto mantém a avaliação por quatro razões, que são também o que ela pode
entregar.

\begin{enumerate}
\item \textbf{Um limite para o efeito médio.} Um intervalo de confiança que
exclua melhoras maiores que algumas dezenas de gramas informa que a lei não
produziu ganho grande na média dos municípios. Não informa sobre quem morava ao
lado do bananal.
\item \textbf{O primeiro estágio.} A medida de quanto avião havia, e de quando
parou, tem valor por si e condiciona a leitura de qualquer resultado.
\item \textbf{A fração exposta.} Medi-la permite dizer que efeito por exposta
seria compatível com o intervalo estimado.
\item \textbf{Os outros elos da cadeia.} Substituição de tecnologia,
intoxicações de trabalhadores, produção e emprego são perguntas da avaliação
com poder próprio, a calcular.
\end{enumerate}

A referência para julgar o efeito mínimo detectável é o menor efeito que
justificaria a política. O segundo ensaio o calcula: é o efeito sobre o peso ao
nascer a partir do qual o benefício supera o custo
(Capítulo~\ref{sec:cb}).

A regra de leitura fica fixada: o intervalo de confiança decide. O efeito mínimo
detectável é instrumento de planejamento e não serve para rebaixar, depois, um
intervalo que exclua o zero, nem para promover um que o inclua.
